% !TEX root = ../main.tex
\corrigesChapitre{7}{Équilibre chimique et sens d'évolution}
\corrigesSouples

\begin{blocLivre}{cor-7-1}
\begin{corrige}{1}{Total ou non total ?}
\demarche{Calculer $x_\mathrm{f}$ à partir de la composition finale et $x_{\max}$ à partir du réactif limitant, puis $\tau=x_\mathrm{f}/x_{\max}$.}
\begin{questions}
\item $\tau=\dfrac{x_\mathrm{f}}{x_{\max}}$ : c'est la proportion de l'avancement maximal effectivement atteinte.
\item A : $x_\mathrm{f}=[\ce{H3O+}]_\mathrm{f}V=\qty{1,0e-2}{mol}=x_{\max}$, \res{\tau=1} ; B : $x_\mathrm{f}=10^{-3{,}4}\times1{,}0=\qty{4,0e-4}{mol}$,
\res{\tau=0{,}040} ; C : $\tau=\frac{0{,}33}{0{,}50}=\res{0{,}66}$.
\item A est \res{\text{totale}} ; B et C sont \res{\text{non totales}}.
\item Le système B contient tous les réactifs et tous les produits : de l'acide éthanoïque (en majorité), des ions
éthanoate et des ions oxonium, dans l'eau.
\end{questions}
\end{corrige}
\end{blocLivre}
\begin{blocLivre}{cor-7-2}
\begin{corrige}{2}{Écrire un quotient de réaction}
\demarche{Produits au numérateur, réactifs au dénominateur, chaque concentration élevée à la puissance de son nombre stœchiométrique.}
\begin{questions}
\item A : $Q_\mathrm{r}=\dfrac{[\ce{CH3COO-}][\ce{H3O+}]}{[\ce{CH3COOH}]}$ ; B : $Q_\mathrm{r}=\dfrac{[\ce{Zn^2+}]}{[\ce{Cu^2+}]}$ ;
C : $Q_\mathrm{r}=[\ce{Ag+}][\ce{Cl-}]$ ; D : $Q_\mathrm{r}=\dfrac{[\ce{I-}]^2[\ce{S4O6^2-}]}{[\ce{I2}][\ce{S2O3^2-}]^2}$ ;
E : $Q_\mathrm{r}=\dfrac{[\ce{FeSCN^2+}]}{[\ce{Fe^3+}][\ce{SCN-}]}$, les concentrations étant divisées par $c^\circ$.
\item L'eau, solvant, et les solides sont des corps purs dont la « concentration » ne varie pas : par convention, ils
\res{\text{n'apparaissent pas}} dans le quotient de réaction.
\item $Q_\mathrm{r}=\dfrac{\num{4.0e-4}}{(\num{2.0e-3})^2}=\res{\num{1.0e2}}$.
\end{questions}
\end{corrige}
\end{blocLivre}
\begin{blocLivre}{cor-7-3}
\begin{corrige}{3}{Prévoir le sens d'évolution}
\demarche{Calculer $Q_{\mathrm{r},\mathrm{i}}$ avec les concentrations initiales, puis le comparer à $K$.}
\begin{questions}
\item A : $Q_{\mathrm{r},\mathrm{i}}=\frac{0{,}10\times\num{1.0e-4}}{0{,}10}=\num{1.0e-4}$ ; B : $\frac{\num{1.0e-3}}{(\num{1.0e-3})^2}=\num{1.0e3}$ ;
C : $\frac{0{,}10}{0{,}10}=1{,}0$ ; D : \mbox{$(\num{1.0e-3})^2=\num{1.0e-6}$}.
\item A : $Q_{\mathrm{r},\mathrm{i}}>K$, \res{\text{sens inverse}} ; B : $Q_{\mathrm{r},\mathrm{i}}>K$, \res{\text{sens inverse}} ;
C : $Q_{\mathrm{r},\mathrm{i}}<K$, \res{\text{sens direct}} ; D : $Q_{\mathrm{r},\mathrm{i}}>K$, \res{\text{sens inverse}}.
\item Dans le mélange D, la réaction évolue dans le sens inverse : il se forme un \res{\text{précipité blanc}} de
chlorure d'argent.
\end{questions}
\end{corrige}
\end{blocLivre}
\begin{blocLivre}{cor-7-4}
\begin{corrige}{4}{Ajouter un réactif à l'équilibre}
\demarche{Un renforcement de la couleur signale la formation d'ions \ce{FeSCN^2+}, donc la présence du réactif qui n'a pas été ajouté.}
\begin{questions}
\item Les tubes 2 et 3 deviennent \res{\text{rouge plus foncé}} : il s'est formé des ions \ce{FeSCN^2+} supplémentaires.
\item Pour former des ions \ce{FeSCN^2+} après l'ajout d'ions \ce{Fe^3+}, il fallait des ions \ce{SCN-} : ils étaient
encore présents. De même, l'évolution du tube 3 montre que des ions \ce{Fe^3+} restaient dans le mélange.
\item Les deux réactifs coexistaient avec le produit dans l'état final : la transformation était \res{\text{non totale}}.
\item L'ajout d'ions \ce{Fe^3+} augmente le dénominateur de $Q_\mathrm{r}$ : $Q_\mathrm{r}$ devient inférieur à $K$ et le
système évolue dans le \res{\text{sens direct}}, jusqu'à ce que $Q_\mathrm{r}$ retrouve la valeur $K$.
\end{questions}
\end{corrige}
\end{blocLivre}
\begin{blocLivre}{cor-7-5}
\begin{corrige}{5}{Un équilibre dynamique}
\demarche{Au départ, seuls A et B sont présents : la réaction dans le sens inverse ne peut pas encore se produire.}
\begin{questions}
\item La \res{\text{courbe a}} : sa vitesse est maximale au départ, quand A et B sont abondants, puis diminue à mesure
qu'ils sont consommés. La courbe b part de zéro, car C et D sont absents au départ.
\item Les deux vitesses deviennent égales à partir de \res{\text{150~min environ}}.
\item Les deux réactions continuent de se produire, à la même vitesse : les entités A, B, C et D se transforment sans
cesse les unes en les autres, d'où le qualificatif \res{\text{dynamique}}.
\item \res{\text{Non}} : autant de molécules de C et de D se forment qu'il en disparaît par unité de temps, la
composition reste constante.
\end{questions}
\schema[corrige]{07-ex05}
\end{corrige}
\end{blocLivre}
\begin{blocLivre}{cor-7-6}
\begin{corrige}{6}{Le taux d'avancement dépend de l'état initial}
\demarche{Le pH donne $[\ce{H3O+}]_\text{éq}$, égale à l'avancement volumique final ; $x_{\max}/V=c$.}
\begin{questions}
\item \ce{HCOOH(aq) + H2O(l) <=> HCOO^-(aq) + H3O^+(aq)}. À l'état final : $n(\ce{HCOOH})=cV-x_\mathrm{f}$ et
$n(\ce{HCOO-})=n(\ce{H3O+})=x_\mathrm{f}$ ; $x_{\max}=cV$.
\item $\tau=\dfrac{x_\mathrm{f}}{x_{\max}}=\dfrac{[\ce{H3O+}]_\text{éq}}{c}$. Pour $S_1$, $[\ce{H3O+}]_\text{éq}=10^{-2{,}90}=\qty{1,26e-3}{mol.L^{-1}}$,
d'où $\res{\tau_1=12{,}6~\%}$. Pour $S_2$, $[\ce{H3O+}]_\text{éq}=10^{-3{,}46}=\qty{3,47e-4}{mol.L^{-1}}$, d'où
$\res{\tau_2=34{,}7~\%}$.
\item $Q_{\mathrm{r},\text{éq}}=\dfrac{[\ce{H3O+}]^2}{c-[\ce{H3O+}]}$ : $S_1$ : $\frac{(\num{1.26e-3})^2}{\num{8.74e-3}}\approx\res{\num{1.8e-4}}$ ;
$S_2$ : $\frac{(\num{3.47e-4})^2}{\num{6.53e-4}}\approx\res{\num{1.8e-4}}$.
\item Le taux d'avancement final \res{\text{dépend de la concentration}} : il augmente avec la dilution. En revanche, le
quotient de réaction à l'équilibre est le même : c'est la constante d'équilibre, $K\approx\num{1.8e-4}$.
\end{questions}
\end{corrige}
\end{blocLivre}
\begin{blocLivre}{cor-7-7}
\begin{corrige}{7}{Une constante indépendante de l'état initial}
\demarche{Convertir $\sigma$ en \unit{S.m^{-1}}, obtenir $[\ce{H3O+}]$ en \unit{mol.m^{-3}}, puis en \unit{mol.L^{-1}}.}
\begin{questions}
\item Les seuls ions présents en quantité notable sont \ce{H3O+} et \ce{CH3COO-}, de même concentration ; la loi de
Kohlrausch donne
\[\sigma=\lambda_{\ce{H3O+}}[\ce{H3O+}]+\lambda_{\ce{CH3COO-}}[\ce{CH3COO-}]=\left(\lambda_{\ce{H3O+}}+\lambda_{\ce{CH3COO-}}\right)[\ce{H3O+}].\]
\item Pour la première solution, $\sigma=160$~$\upmu$S$\cdot$cm$^{-1}$ $=\qty{0,0160}{S.m^{-1}}$, d'où
\[[\ce{H3O+}]_\text{éq}=\frac{0{,}0160}{\num{39.09e-3}}=\qty{0,409}{mol.m^{-3}}=\res{\qty{4,09e-4}{mol.L^{-1}}}.\]
Le tableau rassemble les résultats des quatre solutions (questions 2 à 4).
\begin{center}\small\renewcommand{\arraystretch}{1.2}
\begin{tabular}{@{}l*{4}{c}@{}}\toprule
$c$ ($10^{-3}$ mol$\cdot$L$^{-1}$) & 10,0 & 5,0 & 2,0 & 1,0\\\midrule
$[\ce{H3O+}]_\text{éq}$ ($10^{-4}$ mol$\cdot$L$^{-1}$) & 4,09 & 2,87 & 1,78 & 1,24\\
$Q_{\mathrm{r},\text{éq}}=\dfrac{[\ce{H3O+}]_\text{éq}^2}{c-[\ce{H3O+}]_\text{éq}}$ ($10^{-5}$) & 1,75 & 1,74 & 1,74 & 1,74\\[1.5ex]
$\tau=[\ce{H3O+}]_\text{éq}/c$ (\%) & 4,1 & 5,7 & 8,9 & 12,4\\\bottomrule
\end{tabular}\end{center}
\item Le quotient de réaction à l'équilibre ne dépend pas de la concentration initiale : \res{K\approx\num{1.7e-5}}.
\item Le taux d'avancement final \res{\text{augmente avec la dilution}}, de 4~\% à 12~\% environ.
\end{questions}
\end{corrige}
\end{blocLivre}
\begin{blocLivre}{cor-7-8}
\begin{corrige}{8}{Deux sens d'évolution pour une même réaction}
\demarche{L'argent, solide, n'apparaît pas dans $Q_\mathrm{r}$ ; comparer $Q_{\mathrm{r},\mathrm{i}}$ à $K$ pour chaque système.}
\begin{questions}
\item $Q_\mathrm{r}=\dfrac{[\ce{Fe^2+}][\ce{Ag+}]}{[\ce{Fe^3+}]}$.
\item A : $Q_{\mathrm{r},\mathrm{i}}=\frac{0{,}010\times0{,}010}{0{,}10}=\num{1.0e-3}<K$ : \res{\text{sens direct}}, l'argent est
oxydé, la poudre est en partie consommée. B : $Q_{\mathrm{r},\mathrm{i}}=\frac{0{,}10\times0{,}50}{0{,}010}=5{,}0>K$ :
\res{\text{sens inverse}}, de l'argent se dépose.
\item À l'équilibre, $[\ce{Fe^2+}]=[\ce{Ag+}]=0{,}010+x$ et $[\ce{Fe^3+}]=0{,}10-x$, d'où l'équation proposée, qui équivaut à
$x^2+0{,}33\,x-0{,}0309=0$. La racine positive est $\res{x\approx\qty{0,076}{mol.L^{-1}}}$ : $[\ce{Fe^3+}]_\text{éq}\approx\qty{2,4e-2}{mol.L^{-1}}$
et $\res{\tau\approx76~\%}$.
\item Le taux d'avancement dépend de $K$ mais aussi de l'état initial : ici, les produits sont peu abondants au départ et
le réactif \ce{Fe^3+} est concentré, si bien que le système progresse loin dans le sens direct avant que $Q_\mathrm{r}$
n'atteigne $K$.
\end{questions}
\end{corrige}
\end{blocLivre}
\begin{blocLivre}{cor-7-9}
\begin{corrige}{9}{Le diiode et les ions iodure}
\demarche{Exprimer $Q_\mathrm{r}(x)$ à partir du tableau d'avancement : la fonction $f(x)=Q_\mathrm{r}(x)-K$ s'annule à l'équilibre.}
\begin{questions}
\item À l'équilibre, $[\ce{I2}]=c_1-x$, $[\ce{I-}]=c_2-x$ et $[\ce{I3-}]=x$, d'où $Q_\mathrm{r}=\dfrac{x}{(c_1-x)(c_2-x)}$.
\item \code{return x / ((c1 - x) * (c2 - x)) - K}.
\item On partage l'intervalle $[a\,;\,b]$ en deux : si $f(m)<0$, la solution est entre $m$ et $b$, sinon entre $a$ et $m$.
L'intervalle est divisé par deux à chaque tour, jusqu'à la précision voulue. Il faut que $f$ change de signe sur
$[a\,;\,b]$ pour que la solution s'y trouve.
\item Avec \qty{1,0e-2}{mol.L^{-1}} d'ions iodure, \res{87~\%} du diiode est transformé en ions triiodure ; avec un plus
grand excès d'ions iodure, \res{99~\%} : la solution contient alors presque tout le diiode sous forme d'ions \ce{I3-}.
\end{questions}
\end{corrige}
\end{blocLivre}
\begin{blocLivre}{cor-7-10}
\begin{corrige}{10}{L'estérification, une transformation limitée}
\demarche{Le volume $V$ du mélange s'élimine du quotient de réaction ; écrire $Q_{\mathrm{r},\text{éq}}=K$ en fonction de $x$.}
\begin{questions}
\item Chaque concentration vaut $n/V$ : $Q_\mathrm{r}=\dfrac{\frac{n_\text{ester}}{V}\times\frac{n_\text{eau}}{V}}{\frac{n_\text{acide}}{V}\times\frac{n_\text{alcool}}{V}}$ ;
les volumes se simplifient.
\item À l'équilibre, $\frac{x^2}{(1-x)^2}=4{,}0$, soit $\frac{x}{1-x}=2{,}0$ (racine positive) : $\res{x=\qty{0,67}{mol}}$ et
$\res{\tau=67~\%}$.
\item $\frac{x^2}{(1-x)(2-x)}=4{,}0$ donne $x^2=4(2-3x+x^2)$, soit $3x^2-12x+8=0$. La racine inférieure à 1 vaut
$x=\frac{12-\sqrt{48}}{6}\approx\qty{0,85}{mol}$ : \res{\tau=85~\%}.
\item Éliminer l'eau fait diminuer $Q_\mathrm{r}$, qui devient inférieur à $K$ : le système évolue dans le sens direct.
On augmente ainsi la quantité d'ester obtenue, jusqu'à rendre la transformation presque totale.
\end{questions}
\end{corrige}
\end{blocLivre}
\begin{blocLivre}{cor-7-11}
\begin{corrige}{11}{Une transformation totale}
\demarche{Avec un réactif solide en défaut, l'équilibre ne peut être atteint : la transformation s'arrête quand le solide a disparu.}
\begin{questions}
\item $Q_\mathrm{r}=\dfrac{[\ce{Zn^2+}]}{[\ce{Cu^2+}]}$.
\item $[\ce{Cu^2+}]_\text{éq}=\frac{0{,}10}{\num{1.9e37}}\approx\qty{5e-39}{mol.L^{-1}}$, soit
$\num{5e-39}\times0{,}100\times\num{6.02e23}\approx\res{\num{3e-16}}$ ion : il ne reste en pratique \res{\text{aucun ion}}
\ce{Cu^2+} ; la transformation est totale.
\item $n(\ce{Zn})=\frac{0{,}50}{65{,}4}=\qty{7,6e-3}{mol}$, inférieure à $n(\ce{Cu^2+})=\qty{1,0e-2}{mol}$ : le zinc est limitant. En fin
de transformation, $[\ce{Zn^2+}]=\qty{0,076}{mol.L^{-1}}$ et $[\ce{Cu^2+}]=\qty{0,024}{mol.L^{-1}}$, d'où
$Q_\mathrm{r}\approx3{,}2\ll K$. L'équilibre \res{\text{n'est pas atteint}} : la transformation s'arrête parce que le zinc a
entièrement disparu ; elle est totale.
\item Pour $K>10^4$, à l'équilibre, la concentration d'un des réactifs devient négligeable devant celle des produits :
le taux d'avancement final est pratiquement égal à 1.
\end{questions}
\end{corrige}
\end{blocLivre}
\begin{blocLivre}{cor-7-12}
\begin{corrige}{12}{Mesurer une constante d'équilibre}
\demarche{L'absorbance donne $[\ce{FeSCN^2+}]_\text{éq}=x$ ; le tableau d'avancement donne les concentrations des réactifs.}
\begin{questions}
\item Seuls les ions \ce{FeSCN^2+} absorbent à 470~nm : l'absorbance mesure leur concentration, sans interférence.
Loi de Beer-Lambert : $A=\varepsilon\,\ell\,c$.
\item $[\ce{FeSCN^2+}]_\text{éq}=\frac{A}{\varepsilon\ell}=\frac{0{,}394}{\num{4.7e3}\times1{,}00}=\res{\qty{8,4e-5}{mol.L^{-1}}}$. Le tableau
d'avancement donne alors
\[[\ce{Fe^3+}]_\text{éq}=[\ce{SCN-}]_\text{éq}=\num{1.0e-3}-\num{8.4e-5}=\res{\qty{9,2e-4}{mol.L^{-1}}}.\]
\item Pour chaque mélange, on calcule \mbox{$x=[\ce{FeSCN^2+}]_\text{éq}=A/(\varepsilon\ell)$}, puis
\mbox{$[\ce{Fe^3+}]_\text{éq}=[\ce{Fe^3+}]_\mathrm{i}-x$}, \mbox{$[\ce{SCN-}]_\text{éq}=[\ce{SCN-}]_\mathrm{i}-x$} et
$Q_{\mathrm{r},\text{éq}}=\dfrac{[\ce{FeSCN^2+}]_\text{éq}}{[\ce{Fe^3+}]_\text{éq}\,[\ce{SCN-}]_\text{éq}}$.
\begin{center}\small\renewcommand{\arraystretch}{1.2}
\begin{tabular}{@{}l*{3}{c}@{}}\toprule
 & mélange 1 & mélange 2 & mélange 3\\\midrule
$x$ ($10^{-4}$ mol$\cdot$L$^{-1}$) & 0,838 & 1,56 & 2,17\\
$[\ce{Fe^3+}]_\text{éq}$ ($10^{-4}$ mol$\cdot$L$^{-1}$) & 9,16 & 18,4 & 7,83\\
$[\ce{SCN-}]_\text{éq}$ ($10^{-4}$ mol$\cdot$L$^{-1}$) & 9,16 & 8,44 & 27,8\\
$Q_{\mathrm{r},\text{éq}}$ & 99,9 & 100 & 99,6\\\bottomrule
\end{tabular}\end{center}
Les trois valeurs sont égales, aux erreurs de mesure près : \res{K\approx\num{1.0e2}}.
\item Mélange 1 : $\tau=\frac{\num{8.4e-5}}{\num{1.0e-3}}\approx\res{8~\%}$ ; mélange 3 : $\tau=\frac{\num{2.17e-4}}{\num{1.0e-3}}\approx\res{22~\%}$.
Un excès d'ions thiocyanate augmente la proportion d'ions fer(III) transformés.
\end{questions}
\end{corrige}
\end{blocLivre}
\begin{blocLivre}{cor-7-13}
\begin{corrige}{13}{Dissoudre un précipité de chlorure d'argent}
\demarche{Le solide n'apparaît pas dans $Q_\mathrm{r}$ ; exprimer les concentrations à l'équilibre en fonction de la quantité dissoute par litre.}
\begin{questions}
\item $Q_{\mathrm{r},1}=[\ce{Ag+}][\ce{Cl-}]$ et $Q_{\mathrm{r},2}=\dfrac{[\ce{Ag(NH3)2+}][\ce{Cl-}]}{[\ce{NH3}]^2}$.
\item Dans l'eau pure, $[\ce{Ag+}]=[\ce{Cl-}]=s$ et $s^2=K_1$ : $s=\sqrt{\num{1.8e-10}}=\res{\qty{1,3e-5}{mol.L^{-1}}}$, soit
$\num{1.3e-5}\times143{,}3\approx\res{\qty{1,9e-3}{g.L^{-1}}}$.
\item Au départ, il n'y a ni ions \ce{Ag(NH3)2+} ni ions \ce{Cl-} : $Q_{\mathrm{r},\mathrm{i}}=0<K_2$, le système évolue dans le
sens direct, \res{\text{le solide se dissout}}.
\item Si $s'$ est la quantité dissoute par litre, à l'équilibre $[\ce{Ag(NH3)2+}]=[\ce{Cl-}]=s'$ et $[\ce{NH3}]=1{,}0-2s'$ :
\[\frac{s'^2}{(1{,}0-2s')^2}=K_2,\quad\text{soit}\quad\frac{s'}{1{,}0-2s'}=\sqrt{K_2}=0{,}054\quad\text{et}\quad s'\approx\qty{0,049}{mol.L^{-1}},\]
c'est-à-dire $\res{\qty{7,0}{g}}$ par litre, environ \res{\num{3600}} fois plus que dans l'eau pure.
\end{questions}
\end{corrige}
\end{blocLivre}
\begin{blocLivre}{cor-7-14}
\begin{corrige}{14}{Extraire le diiode d'une solution aqueuse}
\demarche{Écrire la conservation de la quantité de diiode et la relation d'équilibre entre les deux phases.}
\begin{questions}
\item Le cyclohexane, non miscible à l'eau, est \res{\text{moins dense}} qu'elle (densité 0,78) : il surnage.
\item La conservation du diiode s'écrit $n_0=n_\mathrm{a}+n_\mathrm{o}$ et l'équilibre de partage
$n_\mathrm{o}/V_\mathrm{o}=K\,n_\mathrm{a}/V_\mathrm{a}$, soit $n_\mathrm{o}=n_\mathrm{a}\,K\,V_\mathrm{o}/V_\mathrm{a}$. Ainsi
\[\frac{n_\mathrm{a}}{n_0}=\frac{1}{1+K\,V_\mathrm{o}/V_\mathrm{a}}=\frac{1}{1+50\times0{,}20}=\frac{1}{11}\approx\res{9{,}1~\%}.\]
Alors $[\ce{I2}]_\text{aq}=\qty{9,1e-5}{mol.L^{-1}}$ et $[\ce{I2}]_\text{cyclo}=K\,[\ce{I2}]_\text{aq}=\qty{4,5e-3}{mol.L^{-1}}$.
\item Avec deux extractions de 10~mL, il reste $\left(\frac{1}{1+50\times0{,}10}\right)^2=\frac{1}{36}\approx\res{2{,}8~\%}$ du diiode dans
l'eau, contre 9,1~\% avec une seule extraction : \res{\text{deux extractions sont plus efficaces}}.
\item Le diiode est une molécule apolaire ; il est bien plus soluble dans un solvant apolaire comme le cyclohexane que
dans l'eau, solvant polaire.
\end{questions}
\end{corrige}
\end{blocLivre}
\begin{blocLivre}{cor-7-15}
\begin{corrige}{15}{La mutarotation du glucose}
\demarche{Écrire que chaque forme est consommée par une réaction et produite par l'autre ; à l'équilibre, les deux vitesses sont égales.}
\begin{questions}
\item Les deux formes sont dans le même volume : $Q_\mathrm{r}=\dfrac{[\beta]}{[\alpha]}=\dfrac{b}{a}$.
\item Pendant \code{dt}, la forme $\alpha$ disparaît à la vitesse $v_1$ et se forme à la vitesse $v_{-1}$ : sa proportion
varie de $(v_{-1}-v_1)\,\mathrm{d}t$, et celle de $\beta$ de l'opposé (méthode d'Euler). Les deux variations se compensent
exactement : \res{a+b\text{ reste égal à 1}}.
\item À l'équilibre, \res{64~\%} de forme $\beta$ et 36~\% de forme $\alpha$, quel que soit l'état initial : les deux
simulations tendent vers la même composition.
\item À l'équilibre, $v_1=v_{-1}$, donc $k_1[\alpha]=k_{-1}[\beta]$ et $K=\frac{[\beta]}{[\alpha]}=\frac{k_1}{k_{-1}}=\frac{0{,}016}{0{,}009}\approx\res{1{,}78}$.
Le programme affiche $b/a=1{,}77$ : l'équilibre est presque atteint au bout de 240~min.
\item $K=\frac{k_1}{k_{-1}}$ ne change pas : la composition finale \res{\text{est la même}}. Les vitesses étant plus grandes,
l'équilibre est atteint \res{\text{plus vite}}.
\end{questions}
\end{corrige}
\end{blocLivre}
